\texttt{emcee} \citep{ForemanMackey13} implements the
affine-invariant ensemble sampler of \citet{GoodmanWeare10}, in which
a set of ``walkers'' explore the parameter space collectively, using
the positions of complementary walkers to propose new steps. Because
the proposal distribution adapts to the covariance structure of the
target, the algorithm is insensitive to linear degeneracies among
parameters and requires no hand-tuning of a proposal covariance
matrix. We initialize 100 walkers in a tight ball around the maximum
likelihood estimate and run each for $2\times10^{4}$ steps, discarding
samples until the effective sample size exceeds $10^{4}$.

\texttt{dynesty} \citep{Speagle20} implements dynamic nested sampling
\citep{Skilling04,Higson19}, an algorithm that simultaneously
estimates the Bayesian evidence $\mathcal{Z}$ and the posterior. Live
points are drawn from the prior, and at each iteration the
lowest-likelihood point is replaced by a new sample drawn from within
an iso-likelihood contour. We use multi-ellipsoidal bounding with
random slice sampling \citep{Handley15a,Handley15b} and 50 initial
live points, adding live points dynamically until the estimated
evidence error satisfies $\Delta\ln\mathcal{Z} < 1$.

\texttt{PyMultiNest} \citep{Buchner14} provides a Python interface to
the \texttt{MultiNest} library \citep{Feroz09,Feroz19}, which
implements the multimodal nested sampling algorithm. \texttt{MultiNest}
partitions the prior volume into multiple ellipsoidal clusters,
enabling efficient sampling of posteriors with well-separated modes.
We run with 500 live points and a sampling efficiency of 0.8, and
terminate when the remaining evidence contribution falls below a
tolerance of $\Delta\ln\mathcal{Z} = 0.5$.

\texttt{UltraNest} \citep{Buchner21} implements a
reactively-allocated nested sampling algorithm that adjusts the number
of live points during the run to maintain a target accuracy in both
the evidence and the posterior. The algorithm uses the MLFriends
region \citep{Buchner19} to draw new live points, which provides
robust performance in high-dimensional or multimodal settings without
requiring user-specified bounding geometries. We run with a minimum of
500 live points and terminate at $\Delta\ln\mathcal{Z} < 0.5$.
